\subsection{A.4: both openness assertions are false}

The following construction is the infinite tensor-product cusp example
identified in the literature as Hochster's
\emph{Non-openness of loci in Noetherian rings} (1973), Example~1
\cite{rootHochster1973}.
We give every step needed here independently; the original paper has not
yet been opened.  Its publication details were checked against the
author's bibliography.  Our argument for the Frobenius flat locus below
does not require an appeal to Kunz's theorem.

\begin{lemma}[Proved: an explicit noetherian ring of infinitely many cusps]\label{SA23-A-10}
Let $k$ be any field, let $z_1,z_2,\ldots$ be algebraically independent
over $k$, and set
\[
 D=k[z_i^2,z_i^3\mid i\geq1]\subset k[z_i\mid i\geq1],\quad
 P_i=(z_i^2,z_i^3)D,\quad
 S=D\setminus\bigcup_{i\geq1}P_i,\quad R=S^{-1}D.
\]
Then $R$ is a one-dimensional noetherian domain.  Its maximal ideals
are precisely $\mathfrak m_i=S^{-1}P_i$, and
\[
 R_{\mathfrak m_i}\simeq
 L_i[z_i^2,z_i^3]_{(z_i^2,z_i^3)},\qquad
 L_i=\operatorname{Frac}\bigl(k[z_j^2,z_j^3\mid j\ne i]\bigr).
\]
Its regular locus is exactly $\{(0)\}$ and is not open.
\end{lemma}
\begin{proof}
All rings embed in a polynomial ring, so $D$ and $R$ are domains.
The homomorphism setting $z_i=0$ has kernel $P_i$ and has a domain
as target, so $P_i$ is prime and $S$ is multiplicatively closed.
A nonzero polynomial in $D$ uses only finitely many variables; for
every other index it does not belong to $P_i$.  Thus a nonzero
element belongs to only finitely many of the $P_i$.

We need a prime-avoidance argument for this infinite family.
For a nonzero $f\in D$ put $B(f)=\{i:f\in P_i\}$, a finite set.
If $f,g\in D$ and $j$ occurs in neither polynomial, then
\[
 B(f+z_j^2g)=B(f)\cap B(g).
\]
For $i\ne j$ this follows by taking coefficients of the independent
variable $z_j$ after setting $z_i=0$.  For $i=j$ the image is the
nonzero polynomial $f$, so neither side contains $j$.
Suppose an ideal $I\subset D$ is not contained in any $P_i$ and
contains a nonzero $f$.  For each $i\in B(f)$ choose $g_i\in I$
outside $P_i$.  Applying the identity successively with unused
indices produces an element of $I$ outside every $P_i$.
It lies in $I\cap S$.  Hence every ideal disjoint from $S$ is
contained in some $P_i$.

Each $S^{-1}P_i$ is maximal.  Indeed, every nonzero class modulo
$P_i$ is represented by a nonzero $h$ using only variables other
than $z_i$.  The element $h+z_i^2$ belongs to $S$: modulo $P_i$
it is $h$, and modulo any other $P_j$ it is a polynomial in $z_i$
whose coefficient of $z_i^2$ is $1$.  Its image modulo $P_i$ is
$h$, so $h$ becomes invertible in $R/S^{-1}P_i$.
The prime-avoidance argument now shows that these are all the
maximal ideals.  They are distinct, since $z_i^2\in P_i$ but
$z_i^2\notin P_j$ for $j\ne i$.

Localizing $D$ at $P_i$ inverts all nonzero elements in the
other-variable subring.  It gives exactly the displayed local
cusp ring over $L_i$, which is noetherian and one-dimensional.
To prove that $R$ itself is noetherian, let $I\subset R$ be a
nonzero ideal and choose $0\ne f\in I$.  Only finitely many maximal
ideals contain $f$.  At each such maximal ideal choose finitely
many elements of $I$ that generate its localization.  Let $J$ be
the ideal generated by these elements and $f$.  Then $I_{\mathfrak m}
=J_{\mathfrak m}$ for every maximal ideal $\mathfrak m$, including
those not containing $f$, where both are the unit ideal.
An $R$-module that is zero at every maximal ideal is zero: a
nonzero element would have a proper annihilator contained in a
maximal ideal and would remain nonzero there.  Apply this to
$I/J$.  Thus $I=J$ is finitely generated.

The local cusp ring has maximal ideal generated minimally by
$z_i^2,z_i^3$: their classes are linearly independent modulo its
square, because no monomial of degree $2$ or $3$ occurs in that
square.  Its embedding dimension is $2$ and its dimension is $1$,
so it is not regular, in every characteristic.  The remaining
local ring is the fraction field, which is regular.  Finally a
nonempty principal open $D_R(f)$ with $f\ne0$ contains all but
finitely many maximal ideals, so $\{(0)\}$ cannot be open.
\end{proof}

\begin{theorem}[Refuted: openness of the formally smooth fibre locus]\label{SA23-A-11}
For a flat map of noetherian rings $A\to B$, the locus of points
whose local ring in the fibre is formally smooth over the residue
field, for its maximal-ideal topology, need not be open.
Explicitly, let $A=\mathbf F_p$ and let $B=R$ in
Lemma~\ref{SA23-A-10} over $k=\mathbf F_p$.  The locus is exactly
$\{(0)\}$.
\end{theorem}
\begin{proof}
The map is flat because every vector space over a field is flat.
Both rings are noetherian.  At the generic point the local ring
is a field extension of the perfect field $\mathbf F_p$, hence
is separable and formally smooth even for the discrete topology.
At every maximal ideal the local ring is a singular cusp, as
proved above.  A noetherian local algebra formally smooth over a
field in its maximal-ideal topology is regular; equivalently over
a perfect field it is geometrically regular.  This is the
field case of the opened geometric-regularity criterion
\cite[Tag 07NQ, Proposition 15.41.5]{rootStacksRegular}.
Thus no maximal ideal belongs to the locus, and the preceding
lemma proves its nonopenness.  The same example also refutes the
assertion if ``formally smooth'' is interpreted with discrete
topologies, since discrete formal smoothness implies its
maximal-ideal-topology version.
\end{proof}

\begin{theorem}[Refuted: openness of the flat locus]\label{SA23-A-12}
Let $R$ be the explicit noetherian domain of
Lemma~\ref{SA23-A-10} over $\mathbf F_p$.  For its absolute
Frobenius homomorphism
$F_R:R\to R$, $a\mapsto a^p$, the flat locus in the target
spectrum is exactly $\{(0)\}$.  It is not open and contains no
nonempty open subset.
\end{theorem}
\begin{proof}
Frobenius has inverse image $\mathfrak q$ on every prime
$\mathfrak q$, because $a^p\in\mathfrak q$ if and only if
$a\in\mathfrak q$.  At the generic point its source and target
are fields, so its target is a vector space and therefore a flat
module over the source.

At $\mathfrak m_i$ put
$T=L_i[z^2,z^3]_{(z^2,z^3)}$, $x=z^2$, $y=z^3$, and
$\mathfrak n=(x,y)$.  Let
$N=\ker(T^2\to T)$ for the $T$-linear map $(a,b)\mapsto ax+by$.
Because $x,y$ form a basis of $\mathfrak n/\mathfrak n^2$ over
$L_i$, every element of $N$ has both entries in $\mathfrak n$.
If the $T$-module $T$ defined by restriction along Frobenius were
flat, tensoring this sequence would say that every relation
between $x^p,y^p$ comes from $N\otimes_{T,F_T}T$.
The image of this module in $T^2$ has both coordinates in
$F_T(\mathfrak n)T\subset\mathfrak n$.
But $z^p\in T$ for every prime $p\geq2$, because every integer
at least $2$ belongs to the semigroup generated by $2,3$.
The vector $(z^p,-1)$ is a relation:
\[
 z^p x^p-y^p=z^{3p}-z^{3p}=0.
\]
Its second coordinate is a unit.  It cannot belong to that image,
which contradicts flatness.  Hence Frobenius is nonflat at every
maximal ideal, completing the exact determination of the locus.
\end{proof}

\begin{proposition}[Known, with hypotheses checked: finite-presentation repair]\label{SA23-A-13}
If $A\to B$ is of finite presentation, its
flat locus in $\operatorname{Spec}B$ is open.  If additionally
$B$ is flat over $A$, the locus of geometrically regular fibre
points is the smooth locus and is open.
\end{proposition}
\begin{proof}
The first assertion is the openness of the flat locus for a
finitely presented module on a finitely presented morphism,
applied to $\mathcal O_{\operatorname{Spec}B}$
\cite[Theorem 37.15.1]{rootStacksOpenFlat}.
For the second assertion, a finite type scheme over a field is
geometrically regular at a point if and only if it is smooth there
\cite[Lemma 33.12.6]{rootStacksSmoothFibres}. The pointwise
finite-presentation criterion then combines that fibre condition
with flatness to give smoothness at the chosen point
\cite[Lemma 10.137.16]{rootStacksSmoothFibres}.
This identifies the two loci; the smooth locus is open by the local
definition of smoothness. Neither
counterexample above is asserted to be of finite presentation.
\end{proof}
